\chapter{Clairvoyance Gap for Symmetric Metrics}\label{ch:sym-clair}

\providecommand{\Gammaone}{\Gamma_{1}}

This chapter gives a sharp result for symmetric stochastic TSP instances with
one stochastic client.  We first construct a balanced three-chain family
with one client active with probability~$1/2$.  Its clairvoyance gap is
\[
  \frac{10m+2}{9m+2}>\frac{11}{10}
\]
for every integer $m\ge3$, and converges to~$10/9$.  We then prove a matching
upper bound for this activation probability.

More generally, for $p\in[0,1]$, let
\[
  \Gammaone(p)
  :=\sup_{\mathcal I}
  \frac{\optadapt(\mathcal I)}{\optpost(\mathcal I)},
\]
where the supremum ranges over symmetric instances with exactly one
stochastic client of activation probability~$p$ and any number of
deterministic clients.  We prove
\[
  \Gammaone(p)\le\frac{2+p}{2+p^2}.
\]
At $p=1/2$, the construction and upper bound coincide, giving
$\Gammaone(1/2)=10/9$.

\section{The balanced three-chain construction}

Fix an integer $m\ge3$ and set
\begin{equation}\label{eq:HUV}
  H:=2m,
  \qquad
  U:=3m+2,
  \qquad
  V:=4m=2H.
\end{equation}
Thus $U<V$ and $2U>V$.

\begin{construction}[Three chains and one stochastic client]\label{con:three-chain}
The underlying undirected graph contains the depot $r=a_0$, one
stochastic client $s$, and three paths
\[
  A=(a_0,a_1,\ldots,a_m),
  \qquad
  B=(b_0,b_1,\ldots,b_m),
  \qquad
  C=(c_0,c_1,\ldots,c_m).
\]
Every path edge has length $2$, so each chain has total length $H=2m$.
Add four \emph{outer edges} of length $U$,
\begin{equation}\label{eq:outer-edges}
  a_0b_m,
  \qquad
  a_0c_m,
  \qquad
  a_mb_0,
  \qquad
  a_mc_0,
\end{equation}
and three \emph{phase edges} of length $V$,
\begin{equation}\label{eq:phase-edges}
  b_mc_0,
  \qquad
  b_0s,
  \qquad
  c_ms.
\end{equation}
The metric $d$ is the shortest-path metric of this graph.  Every
client other than $s$ is deterministic, and $s$ is active with probability
$1/2$.
\end{construction}

Figure~\ref{fig:balanced-three-chain} shows the topology of the construction.
The edge lengths are chosen by balancing the intended tours against the
cheapest competing crossing patterns.

\begin{figure}[htbp]
\centering
\begin{tikzpicture}[
    x=1cm,y=1cm,
    deterministic/.style={rectangle,draw=black,fill=white,minimum width=8.5mm,
      minimum height=7mm,inner sep=0pt,font=\scriptsize},
    depot/.style={deterministic,fill=gray!15},
    stochastic/.style={circle,draw=black,fill=gray!30,minimum size=7.5mm,
      inner sep=0pt,font=\scriptsize},
    internal/.style={black,line width=0.75pt},
    outer/.style={black!70,densely dashed,line width=0.8pt},
    phase/.style={black,dash pattern=on 4pt off 2pt on 1pt off 2pt,
      line width=0.8pt},
    rowlabel/.style={font=\small\bfseries,anchor=east}
]
\node[depot]         (A0) at (0, 2.4) {$a_0$};
\node[deterministic] (A1) at (1.5, 2.4) {$a_1$};
\node[deterministic] (Am) at (6, 2.4) {$a_m$};

\node[deterministic] (Bm) at (0, 0) {$b_m$};
\node[deterministic] (B1) at (4.5, 0) {$b_1$};
\node[deterministic] (B0) at (6, 0) {$b_0$};

\node[deterministic] (C0) at (0,-2.4) {$c_0$};
\node[deterministic] (C1) at (1.5,-2.4) {$c_1$};
\node[deterministic] (Cm) at (6,-2.4) {$c_m$};
\node[stochastic]    (S)  at (7.3,-1.2) {$s$};

\draw[outer] (A0)--(Bm);
\draw[outer] (A0)--(Cm);
\draw[outer] (Am)--(B0);
\draw[outer] (Am)--(C0);

\draw[phase] (Bm)--(C0);
\draw[phase] (B0)--(S);
\draw[phase] (Cm)--(S);

\draw[internal] (A0)--(A1)--(Am);
\draw[internal] (Bm)--(B1)--(B0);
\draw[internal] (C0)--(C1)--(Cm);

\node[fill=white,inner sep=3pt] at (3.45, 2.4) {$\cdots$};
\node[fill=white,inner sep=3pt] at (2.1, 0) {$\cdots$};
\node[fill=white,inner sep=3pt] at (3.45,-2.4) {$\cdots$};

\node[rowlabel] at (-0.65, 2.4) {chain $A$};
\node[rowlabel] at (-0.65, 0) {chain $B$};
\node[rowlabel] at (-0.65,-2.4) {chain $C$};
\node[above=1mm of A0,font=\scriptsize] {depot $r$};
\end{tikzpicture}
\caption{The balanced three-chain graph.  Solid chain edges have length~$2$,
dashed outer edges have length~$U$, and dash-dotted phase edges have
length~$V$.  The gray circle is the only stochastic client.}
\label{fig:balanced-three-chain}
\end{figure}


Let $D$ be the set of deterministic clients.  Consider the deterministic-only
tour
\begin{equation}\label{eq:T0}
\begin{split}
  T_0:\quad
  &a_0,a_1,\ldots,a_m,
  b_0,b_1,\ldots,b_m,\\
  &c_0,c_1,\ldots,c_m,a_0,
\end{split}
\end{equation}
and the all-active tour
\begin{equation}\label{eq:T1}
\begin{split}
  T_1:\quad
  &a_0,a_1,\ldots,a_m,
  c_0,c_1,\ldots,c_m,\\
  &s,b_0,b_1,\ldots,b_m,a_0.
\end{split}
\end{equation}
Their costs are
\begin{align}
  L_0&:=\cost(T_0)=3H+2U+V=16m+4,
  \label{eq:L0}\\
  L_1&:=\cost(T_1)=3H+2U+2V=20m+4.
  \label{eq:L1}
\end{align}
In particular,
\begin{equation}\label{eq:Ldifference}
  L_1-L_0=V.
\end{equation}

\subsection{Expanded walks and chain contributions}

For every ordered pair of metric vertices, fix one shortest path in the
underlying graph, choosing the reverse path for the reverse ordered
pair.  Replacing every metric edge of a tour by its fixed shortest path gives
an \emph{expanded closed walk} of exactly the same cost.  A traversal of one
of the seven non-chain edges in \cref{eq:outer-edges,eq:phase-edges} is called
a \emph{crossing}.  Crossings are counted with multiplicity.

\begin{lemma}[Chain contribution]\label{lem:chain-contribution}
Let $W$ be an expanded closed walk that serves every client of one of the
chains $A,B,C$.
\begin{enumerate}[label=(\roman*)]
  \item The internal edges of that chain contribute at least $H$ to
  $\cost(W)$.
  \item If the chain has exactly two crossing incidences and both are at the
  same endpoint, its internal contribution is at least $2H$.
  \item Suppose all four outer edges in \cref{eq:outer-edges} are traversed
  exactly once.  Then $A$ has two external incidences at each endpoint, and
  its internal contribution is at least
  \[
    2(H-2)=2H-4.
  \]
\end{enumerate}
\end{lemma}

\begin{proof}
A chain is a path whose only external attachments are its two endpoints.  If
two internal edges were unused, the vertices in a component strictly between
two unused edges would be disconnected from the expanded walk.  Hence at most
one internal edge can be unused.

If every internal edge is used, the contribution is at least $H$.  If exactly
one length-$2$ edge is unused, each of the two remaining path components is
attached to the rest of the walk at only one endpoint.  Since the walk is
closed, every used edge in either component is traversed at least twice.
Consequently the contribution is at least $2(H-2)\ge H$, proving (i).

For (ii), both external incidences occur at one endpoint.  To serve the other
endpoint and return, every internal edge must be traversed at least twice, for
a contribution of at least $2H$.

For (iii), the number of external incidences at each endpoint of $A$ is even.
Because the expanded walk has even degree at every vertex, the multiplicity of
every used edge of the path $A$ is even.  At most one internal edge can be
unused, so the contribution is at least $2(H-2)$.
\end{proof}

\subsection{Rigidity of the two posterior tours}

\begin{lemma}[Inactive realization]\label{lem:inactive-rigidity}
When $s$ is inactive, the optimal tour cost is $L_0$.  Moreover, every tour of
cost strictly less than $L_1$ has, up to reversal, exactly the crossing
skeleton
\begin{equation}\label{eq:inactive-skeleton}
  a_mb_0,
  \qquad
  b_mc_0,
  \qquad
  c_ma_0.
\end{equation}
\end{lemma}

\begin{proof}
Consider an expanded closed walk serving all deterministic clients.  If it
uses $s$ as an internal vertex of one of its shortest paths, then it has at
least two crossings incident to $A$, each of cost $U$, and at least two
crossings incident to $s$, each of cost $V$.  By
\cref{lem:chain-contribution}(i), its cost is at least
\[
  3H+2U+2V=L_1.
\]
Thus a tour of cost below $L_1$ does not traverse $s$.

Contract each of $A,B,C$ to one vertex.  Let $x,y,z$ be the respective numbers
of $A$--$B$, $A$--$C$, and $B$--$C$ crossings.  The contracted multigraph is
connected and Eulerian, so
\[
  x+y\equiv x+z\equiv y+z\equiv0\pmod2.
\]
Hence $x,y,z$ have the same parity.

Suppose first that they are odd.  Then $x,y,z\ge1$.  If one is at least $3$,
the crossing cost is at least $4U+V$, and the total cost is at least
\[
  3H+4U+V>L_1,
\]
because $2U>V$.  Therefore a tour below $L_1$ must have
$x=y=z=1$.  Its contracted walk is the triangle through $A,B,C$ and its
crossing cost is $2U+V$.

The $B$--$C$ crossing is necessarily $b_mc_0$.  For $B$ to use opposite
endpoints, the $A$--$B$ crossing must be $a_mb_0$; for $C$ to use opposite
endpoints, the $A$--$C$ crossing must be $a_0c_m$.  These choices also make
$A$ use opposite endpoints.  Conversely, if one chain uses the same endpoint
twice, at least one other chain does as well.  By
\cref{lem:chain-contribution}, a noncanonical port choice has internal
contribution at least $2H+2H+H=5H$.  Its total cost is therefore at least
\[
  5H+2U+V=L_1,
\]
where we used $V=2H$.  Hence every tour below $L_1$ in the odd case has the
skeleton in \cref{eq:inactive-skeleton}, up to reversal.  Its cost is at least
$3H+2U+V=L_0$, and $T_0$ attains this bound.

Now suppose $x,y,z$ are even.  Connectivity implies that at least two of them
are positive.  If $z>0$, then there are at least two $B$--$C$ crossings and at
least two crossings incident to $A$, so the total cost is at least
\[
  3H+2U+2V=L_1.
\]
Assume therefore that $z=0$.  Then $x,y\ge2$.  If $x+y\ge6$, the total cost is
at least $3H+6U>L_1$.  The only remaining possibility is
\[
  x=y=2,
  \qquad
  z=0.
\]
If $B$ or $C$ uses the same endpoint twice, the total internal contribution is
at least $4H$, and the cost is greater than $L_1$.  Otherwise both $B$ and $C$
use opposite endpoints.  Hence each of the four outer edges is used exactly
once.  By \cref{lem:chain-contribution}(iii), the internal contribution is at
least
\[
  H+H+(2H-4)=4H-4.
\]
Thus the total cost is at least
\[
  4H-4+4U
  =20m+4
  =L_1.
\]
This excludes the even case below $L_1$ and completes the proof.
\end{proof}

\begin{lemma}[Active realization]\label{lem:active-rigidity}
When $s$ is active, the optimal tour cost is $L_1$.  Moreover, every tour of
cost strictly less than $L_1+V$ has, up to reversal, exactly the crossing
skeleton
\begin{equation}\label{eq:active-skeleton}
  a_mc_0,
  \qquad
  c_ms,
  \qquad
  sb_0,
  \qquad
  b_ma_0.
\end{equation}
\end{lemma}

\begin{proof}
Any expanded closed walk serving $s$ has at least two crossings incident to
$A$, each of cost $U$, and at least two crossings incident to $s$, each of cost
$V$.  Therefore
\[
  \cost(W)\ge3H+2U+2V=L_1
\]
by \cref{lem:chain-contribution}(i), and $T_1$ attains this value.

Now consider a tour of cost below $L_1+V$, and let $q$ be its total number of
crossings.  After contracting $A,B,C$ and retaining $s$, we obtain a connected
Eulerian multigraph on four vertices.  Hence every contracted vertex has
positive even degree and $q\ge4$.

If $q=4$, every contracted degree is $2$, so the quotient is the four-cycle
\[
  A-B-s-C-A
\]
up to reversal.  The two edges incident to $s$ are necessarily $b_0s$ and
$c_ms$.  In order that $B$ and $C$ use opposite endpoints, the remaining two
crossings must be $a_0b_m$ and $a_mc_0$; these choices also make $A$ use
opposite endpoints.  Any other port choice makes at least two chains use the
same endpoint twice.  Its internal contribution is then at least $5H$, so its
cost is at least
\[
  5H+2U+2V=L_1+V.
\]
Thus strict inequality forces exactly the skeleton in
\cref{eq:active-skeleton}, up to reversal.

If $q=5$, let $q_A$ be the number of crossings incident to $A$ and $q_s$ the
number incident to $s$.  Both are positive and even.  Every crossing is either
incident to $A$, incident to $s$, or is a $B$--$C$ crossing.  Therefore
$q=5$ forces
\[
  q_A=2,
  \qquad
  q_s=2,
  \qquad
  q_{BC}=1.
\]
The crossing cost is at least $2U+3V$, and hence the total cost is at least
$L_1+V$.

Finally, if $q\ge6$, the two mandatory $A$-crossings and two mandatory
$s$-crossings are supplemented by at least two further crossings.  Every
additional crossing costs at least $U$, so
\[
  \cost(W)
  \ge3H+4U+2V
  >3H+2U+3V
  =L_1+V,
\]
where the strict inequality is $2U>V$.  This proves the rigidity statement.
\end{proof}

\subsection{The adaptive optimum}

\begin{lemma}[Prefix incompatibility]\label{lem:prefix-incompatibility}
No adaptive policy can simultaneously have realized cost below $L_1$ when
$s$ is inactive and realized cost below $L_1+V$ when $s$ is active.
\end{lemma}

\begin{proof}
Assume, for contradiction, that such a policy exists.  Expand every movement
using the fixed shortest paths chosen above.  Before the policy calls $s$, all
called clients are deterministic.  Therefore the active and inactive
realizations make exactly the same calls, occupy the same locations, and have
the same expanded-walk prefix up to the moment at which $s$ is called.

By \cref{lem:inactive-rigidity}, the inactive expanded walk uses only the
three crossings in \cref{eq:inactive-skeleton}.  By
\cref{lem:active-rigidity}, the active expanded walk uses only the four
crossings in \cref{eq:active-skeleton}.

In the active skeleton, the component immediately preceding $s$ is either
$C$ or, under reversal, $B$.  That chain has exactly two crossing incidences in
the entire active expanded walk, so all calls to its deterministic clients must
occur in its unique component interval.  Moreover, the crossing by which the
walk enters that interval must occur before the call to $s$.  Indeed, if the
entry crossing occurred only as part of the post-call movement to the active
vertex $s$, none of the clients of the preceding chain could have been called
in that interval before $s$; serving them later would require re-entering the
chain and would create additional crossings.

Consequently, before the call to $s$, the common expanded prefix has already
used either
\[
  a_mc_0
  \qquad\text{or}\qquad
  a_0b_m,
\]
depending on the orientation of the active skeleton.  Neither crossing occurs
in the inactive skeleton \cref{eq:inactive-skeleton}.  This contradicts the
fact that the inactive realization has the same expanded prefix before the
call to $s$.
\end{proof}

\begin{theorem}[Cost of the three-chain instance]\label{thm:three-chain-cost}
For every integer $m\ge3$, the instance in
\cref{con:three-chain} satisfies
\[
  \optpost=18m+4,
  \qquad
  \optadapt=20m+4.
\]
Therefore
\[
  \frac{\optadapt}{\optpost}
  =\frac{20m+4}{18m+4}
  =\frac{10m+2}{9m+2}
  >\frac{11}{10},
\]
and the ratio converges to $10/9$ as $m\to\infty$.
\end{theorem}

\begin{proof}
By \cref{lem:inactive-rigidity,lem:active-rigidity}, the two posterior optimum
costs are $L_0$ and $L_1$.  Since $p_s=1/2$,
\[
  \optpost=\frac{L_0+L_1}{2}=18m+4.
\]

For an adaptive upper bound, follow the order of $T_1$ and call $s$ upon
reaching $c_m$.  If $s$ is active, the realized tour is exactly $T_1$.  If it
is inactive, call $b_0$ next.  The metric distance satisfies
\[
  d(c_m,b_0)\le d(c_m,s)+d(s,b_0)=2V,
\]
which is exactly the cost of the two phase edges used by the active branch.
Thus both realized costs are at most $L_1$, and
$\optadapt\le L_1$.

For the reverse inequality, let $C_0$ and $C_1$ be the inactive and active
realized costs of an arbitrary adaptive policy.  Suppose that
$(C_0+C_1)/2<L_1$.  Posterior optimality gives
$C_0\ge L_0$ and $C_1\ge L_1$, and therefore
\[
  C_0<L_1,
  \qquad
  C_1<2L_1-L_0=L_1+V,
\]
where the last equality follows from \cref{eq:Ldifference}.  This contradicts
\cref{lem:prefix-incompatibility}.  Hence every adaptive policy has expected
cost at least $L_1$, proving $\optadapt=L_1$.

Finally,
\[
  \frac{10m+2}{9m+2}>\frac{11}{10}
  \quad\Longleftrightarrow\quad
  m>2.
\]
\end{proof}

\subsection{Why these edge lengths?}

The choices of $U$ and $V$ are dictated by the two thresholds in the rigidity
argument.  First consider a three-crossing walk that uses the wrong pair of
chain endpoints.  Relative to the intended inactive tour, at least two chains
must each be traversed to an endpoint and back.  Its internal chain
contribution therefore increases from $3H$ to at least $5H$.  Choosing
\[
  V=2H
\]
makes this extra chain cost equal to the difference
$L_1-L_0=V$.  Thus every such wrong-port walk has cost at least~$L_1$, exactly
the threshold needed in \cref{lem:inactive-rigidity}.

The remaining critical competitor in the inactive realization uses all four
outer edges.  By \cref{lem:chain-contribution}(iii), its cost is at least
\[
  4H-4+4U.
\]
We need this value to be at least the cost of the intended active tour,
\[
  L_1=3H+2U+2V.
\]
To make the outer edges as short as possible without allowing this competitor
below~$L_1$, we impose equality:
\[
  4H-4+4U=3H+2U+2V.
\]
Together with $V=2H$, this gives
\[
  U=\frac{3H}{2}+2,
\]
which is exactly $U=3m+2$ when $H=2m$.

If $U$ were smaller, the four-outer-edge competitor would fall below the
threshold~$L_1$ and the rigidity proof would fail.  If $U$ were larger, both
posterior tour costs would increase while their difference~$V$ stayed fixed,
moving the resulting ratio closer to~$1$.  Hence the chosen values are the
natural tight balance for this construction.

\section{A general upper bound with one stochastic vertex}

We now consider an arbitrary symmetric metric instance with exactly one
stochastic client $s$, active with probability $p\in(0,1)$.  Let $D$ be the
set of deterministic clients.  If $D=\varnothing$, then calling $s$ at the
depot is posterior-optimal in both realizations, so the clairvoyance gap is
$1$.  We therefore assume $D\ne\varnothing$.  Let $T_0$ be an optimal tour on
$\{r\}\cup D$, and let $T_1$ be an optimal tour on
$\{r\}\cup D\cup\{s\}$.  Write
\[
  A:=\cost(T_0),
  \qquad
  B:=\cost(T_1),
  \qquad
  \Delta:=B-A.
\]
Shortcutting $s$ from $T_1$ shows that $A\le B$, so $\Delta\ge0$.  The
posterior optimum is
\begin{equation}\label{eq:post-general}
  \optpost=(1-p)A+pB=A+p\Delta.
\end{equation}

We need two local exchange quantities.

First, define the cheapest insertion surcharge of $s$ into $T_0$ by
\begin{equation}\label{eq:alpha}
  \alpha
  :=\min_{uv\in E(T_0)}
  \bigl(d(u,s)+d(s,v)-d(u,v)\bigr).
\end{equation}
Inserting $s$ into the minimizing edge gives a full tour of cost $A+\alpha$,
so
\begin{equation}\label{eq:delta-alpha}
  \Delta\le\alpha.
\end{equation}

Second, let $x,y$ be the two neighbors of $s$ in $T_1$, and define the saving
from shortcutting $s$ by
\begin{equation}\label{eq:beta}
  \beta:=d(x,s)+d(s,y)-d(x,y).
\end{equation}
The triangle inequality gives $\beta\ge0$, and deleting $s$ from $T_1$ gives
a deterministic tour of cost $B-\beta$.  Hence
\begin{equation}\label{eq:beta-delta}
  0\le\beta\le\Delta.
\end{equation}

The key point is that the cheapest insertion surcharge into $T_0$ is
controlled by the deletion saving in $T_1$.

\begin{lemma}[Insertion--deletion exchange]\label{lem:exchange}
The quantities in \cref{eq:alpha,eq:beta} satisfy
\[
  \alpha\le\frac A2+\beta.
\]
\end{lemma}

\begin{proof}
The vertices $x,y$ split the cycle $T_0$ into two $x$--$y$ arcs.  Let $P$ be
the shorter arc.  Then
\[
  \cost(P)\le\frac A2,
  \qquad
  d(x,y)\le\cost(P)\le\frac A2.
\]
By the definition of $\beta$,
\[
  d(x,s)+d(s,y)=d(x,y)+\beta\le\frac A2+\beta.
\]
Therefore one of $x,y$, say $z$, satisfies
\[
  d(z,s)\le\frac A4+\frac\beta2.
\]
Let $z'$ be either neighbor of $z$ in $T_0$.  Inserting $s$ into the edge
$zz'$ costs
\[
\begin{aligned}
  d(z,s)+d(s,z')-d(z,z')
  &\le d(z,s)+\bigl(d(s,z)+d(z,z')\bigr)-d(z,z')\\
  &=2d(z,s)\\
  &\le\frac A2+\beta.
\end{aligned}
\]
Since $\alpha$ is the minimum insertion surcharge, the claim follows.
\end{proof}

\subsection{Two adaptive policies}

\begin{lemma}[Two-policy bound]\label{lem:two-policies}
For every one-stochastic-vertex instance,
\begin{equation}\label{eq:two-policy-bound}
  \optadapt
  \le
  \min\left\{
    A+p\left(\frac A2+\beta\right),
    \;A+\Delta-(1-p)\beta
  \right\}.
\end{equation}
\end{lemma}

\begin{proof}
For the first policy, orient $T_0$ and follow it until reaching one endpoint
$u$ of an edge $uv$ attaining the minimum in \cref{eq:alpha}.  Call $s$ at
$u$.  If $s$ is inactive, continue directly to $v$ and complete $T_0$, at cost
$A$.  If $s$ is active, traverse $u,s,v$ and then complete $T_0$, at cost
$A+\alpha$.  Thus the expected cost is
\[
  A+p\alpha
  \le A+p\left(\frac A2+\beta\right)
\]
by \cref{lem:exchange}.

For the second policy, orient $T_1$ so that the local order is
$x,s,y$.  Follow $T_1$ to $x$ and call $s$.  If $s$ is active, complete
$T_1$, at cost $B$.  If $s$ is inactive, move directly from $x$ to $y$ and
complete the shortcut tour, at cost $B-\beta$.  Its expected cost is
\[
  pB+(1-p)(B-\beta)
  =B-(1-p)\beta
  =A+\Delta-(1-p)\beta.
\]
Taking the better policy gives \cref{eq:two-policy-bound}.
\end{proof}

\subsection{Optimization of the bound}

\begin{theorem}[One stochastic vertex]\label{thm:general-upper}
For every symmetric metric instance with exactly one stochastic client of
activation probability $p\in[0,1]$,
\[
  \frac{\optadapt}{\optpost}
  \le\frac{2+p}{2+p^2}.
\]
Equivalently,
\[
  \Gammaone(p)\le\frac{2+p}{2+p^2}.
\]
\end{theorem}

\begin{proof}
The cases $p=0$ and $p=1$ have no information gap and give ratio $1$, so
assume $p\in(0,1)$.  Because $D\ne\varnothing$ and $d$ is a metric, $A>0$.  Normalize
\[
  \delta:=\frac{\Delta}{A},
  \qquad
  b:=\frac{\beta}{A}.
\]
By \cref{eq:post-general,eq:two-policy-bound},
\begin{equation}\label{eq:normalized-bound}
  \frac{\optadapt}{A}
  \le
  \min\left\{
    1+\frac p2+pb,
    \;1+\delta-(1-p)b
  \right\},
  \qquad
  \frac{\optpost}{A}=1+p\delta.
\end{equation}

Suppose first that $\delta\le p/2$.  The second policy and $b\ge0$ give
\[
  \frac{\optadapt}{A}\le1+\delta.
\]
The function $(1+\delta)/(1+p\delta)$ is increasing in $\delta$ because
$p<1$.  Hence
\[
  \frac{\optadapt}{\optpost}
  \le\frac{1+\delta}{1+p\delta}
  \le\frac{1+p/2}{1+p^2/2}
  =\frac{2+p}{2+p^2}.
\]

Now suppose that $\delta\ge p/2$ and set
\[
  b_0:=\delta-\frac p2\ge0.
\]
If $b\le b_0$, the first term in \cref{eq:normalized-bound} gives
\[
\begin{aligned}
  \frac{\optadapt}{A}
  &\le1+\frac p2+pb\\
  &\le1+\frac p2+p\left(\delta-\frac p2\right)\\
  &=1+p\delta+\frac{p(1-p)}2.
\end{aligned}
\]
If $b\ge b_0$, the second term gives
\[
\begin{aligned}
  \frac{\optadapt}{A}
  &\le1+\delta-(1-p)b\\
  &\le1+\delta-(1-p)\left(\delta-\frac p2\right)\\
  &=1+p\delta+\frac{p(1-p)}2.
\end{aligned}
\]
Thus in either case,
\[
  \frac{\optadapt}{\optpost}
  \le
  1+\frac{p(1-p)/2}{1+p\delta}.
\]
Since $\delta\ge p/2$,
\[
  \frac{\optadapt}{\optpost}
  \le
  1+\frac{p(1-p)/2}{1+p^2/2}
  =\frac{1+p/2}{1+p^2/2}
  =\frac{2+p}{2+p^2}.
\]
This proves the theorem.
\end{proof}

\begin{theorem}[Exact unbiased one-vertex gap]\label{thm:sym-clair}
Among symmetric instances with exactly one stochastic client of probability
$1/2$,
\[
  \Gammaone\!\left(\frac12\right)=\frac{10}{9}.
\]
\end{theorem}

\begin{proof}
By \cref{thm:general-upper},
\[
  \Gammaone\!\left(\frac12\right)
  \le
  \frac{2+1/2}{2+(1/2)^2}
  =\frac{10}{9}.
\]
By \cref{thm:three-chain-cost}, the three-chain family has ratios converging
to $10/9$.  Hence the supremum is exactly $10/9$.
\end{proof}

\begin{remark}[Maximum of the general upper bound]
The function
\[
  f(p):=\frac{2+p}{2+p^2}
\]
is maximized at
\[
  p=\sqrt6-2,
\]
where
\[
  f(p)=\frac{2+\sqrt6}{4}\approx1.112372.
\]
Thus every symmetric instance with one stochastic client, with arbitrary
activation probability, satisfies
\[
  \frac{\optadapt}{\optpost}
  \le\frac{2+\sqrt6}{4}.
\]
The present construction settles the unbiased case exactly; whether the
probability-dependent bound is attainable for biased $p$ is a separate
question.
\end{remark}

\begin{remark}[Extremal parameters of the construction]
For the three-chain family,
\[
  A=L_0=16m+4,
  \qquad
  \Delta=L_1-L_0=4m.
\]
The neighbors of $s$ in $T_1$ are $c_m$ and $b_0$, and
$d(c_m,b_0)=2V$, so shortcutting $s$ saves $\beta=0$.  A direct check of the
three internal-edge families and the three crossing edges of $T_0$ shows that
every insertion surcharge is at least $2V$; equality holds, for example, on
the edge $a_mb_0$.  Thus $\alpha=2V=8m$.  Consequently,
\[
  \frac{\Delta}{A}\longrightarrow\frac14,
  \qquad
  \frac{\beta}{A}=0,
  \qquad
  \frac{\alpha}{A}\longrightarrow\frac12.
\]
These are exactly the limiting equality conditions in the proof of
\cref{thm:general-upper} at $p=1/2$.
\end{remark}
